\documentclass[11pt,a4paper]{article}

\usepackage[utf8]{inputenc}
\usepackage[T1]{fontenc}
\usepackage[french]{babel}
\usepackage{geometry}
\usepackage{xcolor}
\usepackage{hyperref}
\usepackage{tabularx}
\usepackage{longtable}
\usepackage{enumitem}
\usepackage{listings}
\usepackage{titlesec}

\geometry{margin=2.4cm}
\hypersetup{
  colorlinks=true,
  linkcolor=blue!55!black,
  urlcolor=blue!55!black
}

\definecolor{chorusblue}{HTML}{1D4ED8}
\definecolor{chorusgray}{HTML}{F3F4F6}
\definecolor{chorusdark}{HTML}{111827}

\titleformat{\section}
  {\Large\bfseries\color{chorusblue}}
  {\thesection}{0.75em}{}

\titleformat{\subsection}
  {\large\bfseries\color{chorusdark}}
  {\thesubsection}{0.75em}{}

\lstdefinestyle{chorus}{
  backgroundcolor=\color{chorusgray},
  basicstyle=\ttfamily\small,
  breaklines=true,
  frame=single,
  rulecolor=\color{black!15},
  columns=fullflexible
}
\lstset{style=chorus}

\setlist[itemize]{leftmargin=1.2cm}
\setlist[enumerate]{leftmargin=1.2cm}

\title{
  \vspace{-1.5cm}
  \textbf{\Huge Documentation technique}\\
  \vspace{0.3cm}
  \Large Chorus Pilot
}
\author{Projet Chorus}
\date{\today}

\begin{document}

\maketitle
\thispagestyle{empty}

\begin{center}
  \vspace{0.5cm}
  \begin{tabularx}{0.9\textwidth}{|X|}
    \hline
    \textbf{Objet du document} \\
    Ce document regroupe la documentation technique globale du projet Chorus Path
    Finder : architecture, frontend, backend, base de données, intégration Chorus,
    configuration, sécurité, tests, Docker, CI/CD GitHub Actions et exploitation.
    Il sert de support de présentation technique et de référence projet. \\
    \hline
  \end{tabularx}
\end{center}

\newpage
\tableofcontents
\newpage

\section{Présentation technique du projet}

Chorus Pilot est une application web permettant de préparer, contrôler,
synchroniser et soumettre des factures vers Chorus Pro via PISTE.

Le projet combine :

\begin{itemize}
  \item un frontend React, TypeScript et Vite ;
  \item une API backend Python FastAPI ;
  \item une couche métier Python dédiée à Chorus ;
  \item une base PostgreSQL ;
  \item des migrations Alembic ;
  \item une stack Docker de production ;
  \item des workflows GitHub Actions pour la CI, la sécurité et le déploiement.
\end{itemize}

L'objectif technique principal est de séparer clairement l'interface utilisateur
des appels sensibles à Chorus. Le navigateur ne manipule pas directement les
secrets : les appels nécessitant des identifiants techniques passent par le
backend.

\section{Architecture générale}

\subsection{Vue d'ensemble}

\begin{lstlisting}
Utilisateur
  -> Navigateur web
    -> Application React/Vite
      -> nginx
        -> /api/* vers FastAPI
          -> Authentification JWT
          -> Extraction PDF
          -> Services Chorus / PISTE
          -> API DB
            -> PostgreSQL
\end{lstlisting}

\subsection{Responsabilités par couche}

\begin{longtable}{|p{0.25\textwidth}|p{0.68\textwidth}|}
  \hline
  \textbf{Couche} & \textbf{Responsabilités} \\
  \hline
  Frontend &
  Navigation, formulaires métier, extraction côté client lorsque possible,
  préparation des données, affichage des factures, statuts, référentiels et
  paramètres. \\
  \hline
  Backend FastAPI &
  Authentification, endpoints applicatifs, extraction PDF côté serveur,
  communication sécurisée avec Chorus, synchronisation, accès aux données. \\
  \hline
  Module \texttt{chorus\_python} &
  Client Chorus, construction des payloads, authentification Chorus, modèles,
  services métier, annuaire, extraction et accès base de données. \\
  \hline
  PostgreSQL &
  Persistance des factures, référentiels, synchronisations, historiques et
  données applicatives. \\
  \hline
  Docker / nginx &
  Packaging de production, reverse proxy, exposition HTTP et orchestration des
  services. \\
  \hline
  GitHub Actions &
  Contrôles automatiques, tests, sécurité, validation Docker et déploiement
  développement. \\
  \hline
\end{longtable}

\section{Structure du dépôt}

\begin{longtable}{|p{0.33\textwidth}|p{0.60\textwidth}|}
  \hline
  \textbf{Chemin} & \textbf{Rôle} \\
  \hline
  \texttt{src/} &
  Code frontend React, pages, composants, hooks, types et logique Chorus côté
  client. \\
  \hline
  \texttt{src/pages/} &
  Ecrans principaux : tableau de bord, factures, création, détail, références,
  paramètres et connexion. \\
  \hline
  \texttt{src/lib/chorus/} &
  Préparation, normalisation, validation, hydratation et soumission des données
  Chorus côté frontend. \\
  \hline
  \texttt{chorus\_python/} &
  Modules Python métier : client Chorus, service de soumission, extraction,
  annuaire, authentification, configuration et base de données. \\
  \hline
  \texttt{chorus\_python/db/} &
  Entités, schémas, repositories, session, synchronisation et API DB. \\
  \hline
  \texttt{scripts/chorus\_submit\_server.py} &
  Point d'entrée du serveur FastAPI. \\
  \hline
  \texttt{alembic/} &
  Migrations de base de données PostgreSQL. \\
  \hline
  \texttt{public/test-data/} &
  Données JSON de démonstration et de fallback. \\
  \hline
  \texttt{tests/} &
  Tests Python. \\
  \hline
  \texttt{src/test/} &
  Tests TypeScript avec Vitest. \\
  \hline
  \texttt{docker/} &
  Configuration nginx. \\
  \hline
  \texttt{.github/workflows/} &
  Workflows GitHub Actions. \\
  \hline
\end{longtable}

\section{Frontend}

\subsection{Technologies}

Le frontend est basé sur :

\begin{itemize}
  \item React 18 ;
  \item TypeScript ;
  \item Vite ;
  \item React Router ;
  \item TanStack Query ;
  \item Tailwind CSS ;
  \item composants UI Radix ;
  \item Recharts pour les graphiques ;
  \item Vitest pour les tests.
\end{itemize}

\subsection{Routes principales}

\begin{longtable}{|p{0.28\textwidth}|p{0.65\textwidth}|}
  \hline
  \textbf{Route} & \textbf{Ecran} \\
  \hline
  \texttt{/login} &
  Connexion utilisateur. \\
  \hline
  \texttt{/} &
  Tableau de bord. \\
  \hline
  \texttt{/invoices} &
  Liste des factures. \\
  \hline
  \texttt{/invoices/new} &
  Création et soumission d'une nouvelle facture. \\
  \hline
  \texttt{/invoices/:id} &
  Détail d'une facture locale. \\
  \hline
  \texttt{/invoices/chorus/:id} &
  Détail d'une facture issue de Chorus. \\
  \hline
  \texttt{/references/customers} &
  Référentiel clients publics. \\
  \hline
  \texttt{/references/services} &
  Référentiel services destinataires. \\
  \hline
  \texttt{/settings} &
  Paramètres Chorus et connecteur. \\
  \hline
\end{longtable}

\subsection{Gestion de session côté frontend}

Les routes métier sont protégées par un composant de garde. Si la session
utilisateur n'est pas valide, l'utilisateur est redirigé vers l'écran de
connexion. Le frontend charge ensuite la configuration Chorus locale et initialise
les catalogues de références lorsque l'utilisateur est authentifié.

\subsection{Gestion des données}

TanStack Query est utilisé pour :

\begin{itemize}
  \item charger les factures ;
  \item charger les fournisseurs ;
  \item charger les clients publics et services ;
  \item déclencher les synchronisations ;
  \item éviter les rechargements inutiles grâce au cache.
\end{itemize}

\section{Backend FastAPI}

\subsection{Point d'entrée}

Le serveur principal est défini dans :

\begin{lstlisting}
scripts/chorus_submit_server.py
\end{lstlisting}

Il démarre une application FastAPI exposant les endpoints nécessaires au
frontend. Il charge aussi les routes d'authentification et, lorsque la base est
disponible, les routes \texttt{/api/db/*}.

\subsection{Endpoints principaux}

\begin{longtable}{|p{0.35\textwidth}|p{0.58\textwidth}|}
  \hline
  \textbf{Endpoint} & \textbf{Rôle} \\
  \hline
  \texttt{GET /health} &
  Vérification de santé de l'API. \\
  \hline
  \texttt{POST /api/auth/login} &
  Connexion utilisateur. \\
  \hline
  \texttt{GET /api/auth/me} &
  Lecture de la session courante. \\
  \hline
  \texttt{POST /api/chorus/oauth-token} &
  Récupération d'un token OAuth2 Chorus côté serveur. \\
  \hline
  \texttt{POST /api/extract-pdf-text} &
  Extraction de texte depuis un PDF encodé en base64. \\
  \hline
  \texttt{POST /api/chorus/submit} &
  Upload PDF puis soumission de la facture à Chorus. \\
  \hline
  \texttt{POST /api/chorus/annuaire/sync} &
  Synchronisation de l'annuaire Chorus. \\
  \hline
  \texttt{POST /api/chorus/siret/info} &
  Récupération d'informations liées à une structure SIRET. \\
  \hline
  \texttt{GET /api/db/*} &
  Accès aux données applicatives persistées. \\
  \hline
\end{longtable}

\subsection{Rôle du backend}

Le backend centralise les traitements sensibles :

\begin{itemize}
  \item lecture des secrets serveur ;
  \item appel à l'API Chorus ;
  \item obtention des tokens OAuth2 ;
  \item upload des pièces jointes ;
  \item soumission des factures ;
  \item extraction PDF plus fiable via les bibliothèques Python ;
  \item synchronisation des référentiels ;
  \item exposition contrôlée des données au frontend.
\end{itemize}

\section{Intégration Chorus Pro}

\subsection{Flux de soumission}

Le flux technique de soumission est le suivant :

\begin{enumerate}
  \item l'utilisateur importe ou prépare une facture dans le frontend ;
  \item les données sont normalisées et validées ;
  \item le frontend construit un payload de soumission ;
  \item la demande est envoyée au backend ;
  \item le backend récupère un token OAuth2 Chorus ;
  \item le backend charge le PDF vers Chorus ;
  \item Chorus renvoie une référence de pièce jointe ;
  \item le backend soumet la facture en associant cette pièce jointe ;
  \item le retour Chorus est renvoyé au frontend ;
  \item l'interface affiche le résultat à l'utilisateur.
\end{enumerate}

\subsection{Modes de dépôt}

Les modes principaux pris en compte sont :

\begin{itemize}
  \item \texttt{DEPOT\_PDF\_API} : dépôt PDF via API ;
  \item \texttt{SAISIE\_API} : soumission structurée sans PDF principal ;
  \item \texttt{DEPOT\_PDF\_SIGNE\_API} : dépôt PDF signé si le flux est activé.
\end{itemize}

\subsection{Préparation des payloads}

Le frontend contient des modules dédiés à la construction et à la validation des
payloads Chorus. Les données extraites d'un PDF sont transformées en brouillon de
facture, puis enrichies avec les référentiels et le contexte technique Chorus.

Les validations portent notamment sur :

\begin{itemize}
  \item les montants ;
  \item les dates ;
  \item le fournisseur ;
  \item le destinataire ;
  \item le service exécutant ;
  \item le mode de dépôt ;
  \item les pièces jointes ;
  \item les références de paiement, commande ou marché.
\end{itemize}

\section{Extraction PDF}

\subsection{Principe}

L'application tente d'extraire automatiquement le texte des factures PDF afin de
préremplir les champs métier. L'extraction peut se faire côté frontend, mais le
backend Python fournit une extraction plus robuste grâce aux bibliothèques
spécialisées.

\subsection{Nettoyage du texte}

Le projet inclut une étape de nettoyage pour retirer le bruit technique souvent
présent dans les PDF :

\begin{itemize}
  \item en-têtes PDF ;
  \item objets internes ;
  \item marqueurs de flux ;
  \item métadonnées ;
  \item caractères de contrôle.
\end{itemize}

\section{Base de données}

\subsection{Technologies}

La persistance repose sur PostgreSQL. Les migrations sont gérées avec Alembic.
Le dossier \texttt{chorus\_python/db/} contient les entités, schémas,
repositories, sessions et traitements de synchronisation.

\subsection{Données gérées}

La base peut contenir :

\begin{itemize}
  \item factures ;
  \item lignes de facture ;
  \item pièces jointes ;
  \item historiques de statuts ;
  \item fournisseurs ;
  \item clients publics ;
  \item services destinataires ;
  \item logs d'échanges API ;
  \item traces d'audit ;
  \item jobs asynchrones ou synchronisations.
\end{itemize}

\subsection{Migrations}

Les migrations sont stockées dans \texttt{alembic/versions/}. Elles permettent de
faire évoluer le schéma PostgreSQL de manière contrôlée, notamment lors du
démarrage de la stack Docker.

\section{Configuration}

\subsection{Fichiers de configuration}

\begin{longtable}{|p{0.30\textwidth}|p{0.63\textwidth}|}
  \hline
  \textbf{Fichier} & \textbf{Usage} \\
  \hline
  \texttt{.env.example} &
  Exemple de configuration versionné. \\
  \hline
  \texttt{.env} &
  Variables lues par Docker Compose, par exemple le port HTTP exposé. \\
  \hline
  \texttt{.env.local} &
  Secrets et configuration serveur non versionnés. \\
  \hline
  \texttt{vite.config.ts} &
  Configuration Vite. \\
  \hline
  \texttt{docker-compose.prod.yml} &
  Orchestration Docker de production. \\
  \hline
\end{longtable}

\subsection{Variables Chorus}

Les secrets Chorus doivent rester côté serveur. Les variables principales sont :

\begin{lstlisting}
CHORUS_ENV=prod
CHORUS_PROD_CLIENT_ID=
CHORUS_PROD_CLIENT_SECRET=
CHORUS_PROD_TECHNICAL_LOGIN=
CHORUS_PROD_TECHNICAL_PASSWORD=
CHORUS_PROD_SUPPLIER_ID=
CHORUS_PROD_DESTINATAIRE_CODE=
CHORUS_PROD_SERVICE_EXECUTANT_CODE=
CHORUS_PROD_BANK_ACCOUNT_CODE=
\end{lstlisting}

Les variantes de qualification utilisent le préfixe :

\begin{lstlisting}
CHORUS_QUALIF_*
\end{lstlisting}

\subsection{Variables frontend}

Les variables préfixées par \texttt{VITE\_} sont injectées dans le bundle
navigateur. Elles ne doivent pas contenir de secrets de production.

\section{Sécurité}

\subsection{Principes appliqués}

\begin{itemize}
  \item Les secrets Chorus sont lus côté serveur.
  \item Le navigateur passe par le backend pour les appels sensibles.
  \item L'authentification applicative repose sur des tokens JWT.
  \item Les fichiers \texttt{.env} ne doivent pas être versionnés.
  \item Les workflows GitHub bloquent les secrets évidents commis par erreur.
  \item Les dépendances npm et Python sont auditées.
\end{itemize}

\subsection{Limite actuelle}

L'application utilise actuellement un compte applicatif partagé. Une gestion
multi-utilisateurs complète nécessiterait :

\begin{itemize}
  \item une table utilisateurs ;
  \item une gestion des rôles ;
  \item une traçabilité par utilisateur réel ;
  \item des permissions par action ou périmètre.
\end{itemize}

\section{Docker et déploiement}

\subsection{Services Docker}

La stack Docker de production contient :

\begin{itemize}
  \item un service frontend servi par nginx ;
  \item un service backend FastAPI ;
  \item un service PostgreSQL ;
  \item un service de migration Alembic ;
  \item des volumes persistants pour la base et l'annuaire.
\end{itemize}

\subsection{Commande de démarrage}

\begin{lstlisting}
docker-compose -f docker-compose.prod.yml up -d --build
\end{lstlisting}

\subsection{Accès réseau}

L'application peut être exposée :

\begin{itemize}
  \item sur \texttt{localhost} en local ;
  \item sur une adresse IP interne ;
  \item sur un nom DNS interne comme \texttt{chorusPilot.gigalis.local}.
\end{itemize}

Le port est piloté par la variable :

\begin{lstlisting}
APP_HTTP_PORT=8081
\end{lstlisting}

\section{GitHub Actions et CI/CD}

\subsection{Vue d'ensemble}

Le dossier \texttt{.github/workflows/} contient trois workflows :

\begin{itemize}
  \item \texttt{ci.yml} ;
  \item \texttt{security.yml} ;
  \item \texttt{deploy-development.yml}.
\end{itemize}

Le fichier \texttt{.github/dependabot.yml} configure la surveillance automatique
des dépendances.

\subsection{Workflow CI}

Le workflow \textbf{CI} se déclenche sur les branches :

\begin{itemize}
  \item \texttt{main} ;
  \item \texttt{PRE\_PROD\_V0} ;
  \item \texttt{developpement}.
\end{itemize}

Il contient trois jobs :

\begin{longtable}{|p{0.30\textwidth}|p{0.62\textwidth}|}
  \hline
  \textbf{Job} & \textbf{Contrôles} \\
  \hline
  Frontend checks &
  Installation npm, lint, tests Vitest et build Vite. \\
  \hline
  Backend checks &
  Installation Python, compilation des modules et tests backend. \\
  \hline
  Docker Compose config &
  Validation du fichier \texttt{docker-compose.prod.yml}. \\
  \hline
\end{longtable}

Commandes principales :

\begin{lstlisting}
npm ci
npm run lint
npm run test
npm run build
python -m compileall chorus_python scripts
python -m unittest discover -s tests
docker compose --env-file .env.local -f docker-compose.prod.yml config
\end{lstlisting}

\subsection{Workflow Security}

Le workflow \textbf{Security} exécute :

\begin{itemize}
  \item un contrôle des fichiers sensibles suivis par Git ;
  \item une recherche de secrets évidents ;
  \item un audit npm avec \texttt{npm audit --audit-level=high} ;
  \item un audit Python avec \texttt{pip-audit}.
\end{itemize}

Il est déclenché sur push, pull request, lancement manuel et planification
hebdomadaire.

\subsection{Workflow Deploy Development}

Le workflow \textbf{Deploy Development} déploie automatiquement la branche
\texttt{developpement} sur le serveur de développement après une CI réussie.

Il utilise les secrets GitHub suivants :

\begin{itemize}
  \item \texttt{DEV\_SSH\_HOST} ;
  \item \texttt{DEV\_SSH\_USER} ;
  \item \texttt{DEV\_SSH\_PRIVATE\_KEY}.
\end{itemize}

Sur le serveur, il exécute les actions suivantes :

\begin{enumerate}
  \item récupération des derniers commits ;
  \item bascule sur la branche \texttt{developpement} ;
  \item \texttt{git pull --ff-only} ;
  \item reconstruction et relance Docker Compose ;
  \item affichage de l'état des conteneurs.
\end{enumerate}

\subsection{Dependabot}

Dependabot vérifie chaque semaine :

\begin{itemize}
  \item les dépendances npm ;
  \item les dépendances pip ;
  \item les actions GitHub.
\end{itemize}

Il peut ouvrir jusqu'à cinq pull requests par écosystème.

\section{Tests}

\subsection{Tests frontend}

Les tests frontend sont exécutés avec Vitest. Ils couvrent notamment :

\begin{itemize}
  \item les mappers Chorus ;
  \item la normalisation des données ;
  \item la préparation des payloads ;
  \item l'extraction ;
  \item le stockage ;
  \item l'hydratation de soumission ;
  \item les routes PISTE ;
  \item les clients backend.
\end{itemize}

Commande :

\begin{lstlisting}
npm run test
\end{lstlisting}

\subsection{Tests backend}

Les tests Python sont situés dans \texttt{tests/}. Ils couvrent notamment la
configuration, l'annuaire et la soumission de facture.

Commande utilisée par la CI :

\begin{lstlisting}
python -m unittest discover -s tests
\end{lstlisting}

\subsection{Validation Docker}

La CI valide également que le fichier Compose de production est lisible et
cohérent :

\begin{lstlisting}
docker compose --env-file .env.local -f docker-compose.prod.yml config
\end{lstlisting}

\section{Exploitation}

\subsection{Commandes utiles}

\begin{lstlisting}
npm install
npm run dev
npm run build
npm run lint
npm run test
python scripts/chorus_submit_server.py
docker-compose -f docker-compose.prod.yml up -d --build
docker-compose -f docker-compose.prod.yml ps
docker-compose -f docker-compose.prod.yml logs -f chorus-api
\end{lstlisting}

\subsection{Points de contrôle en production}

\begin{itemize}
  \item vérifier que le backend répond sur \texttt{/health} ;
  \item vérifier que le frontend est accessible ;
  \item vérifier que PostgreSQL est démarré ;
  \item vérifier que les migrations sont passées ;
  \item vérifier que \texttt{.env.local} existe sur le serveur ;
  \item tester la connexion Chorus depuis l'écran Paramètres ;
  \item surveiller les logs backend en cas d'échec de soumission.
\end{itemize}

\section{Synthèse des apports techniques}

Le projet met en place une base technique complète :

\begin{itemize}
  \item application web React structurée par pages et modules métier ;
  \item backend FastAPI servant de passerelle sécurisée vers Chorus ;
  \item extraction PDF côté serveur ;
  \item intégration des workflows Chorus Pro via PISTE ;
  \item base PostgreSQL avec migrations ;
  \item déploiement Docker reproductible ;
  \item validation automatique par GitHub Actions ;
  \item audit de sécurité et dépendances ;
  \item déploiement automatique de la branche de développement ;
  \item documentation utilisateur, projet et technique.
\end{itemize}

\section{Axes d'amélioration}

Les évolutions techniques possibles sont :

\begin{itemize}
  \item gestion multi-utilisateurs complète ;
  \item rôles et permissions ;
  \item audit applicatif par utilisateur réel ;
  \item pipeline de déploiement préproduction et production ;
  \item scans de sécurité plus avancés ;
  \item supervision applicative ;
  \item alertes en cas d'échec de synchronisation ou de rejet Chorus ;
  \item sauvegardes et restauration documentées pour PostgreSQL ;
  \item génération automatique d'une documentation PDF depuis la CI.
\end{itemize}

\section{Conclusion}

Chorus Pilot dispose d'une architecture technique cohérente pour traiter
les factures Chorus Pro de bout en bout. Le projet couvre l'interface métier, les
traitements backend, l'intégration Chorus, la persistance, le déploiement Docker,
les tests et l'automatisation GitHub. Cette base permet une exploitation interne
fiable et prépare les futures évolutions vers une plateforme plus complète,
traçable et industrialisée.

\end{document}
